
%%%%%%%%%%%%%%%%%%%%%%%%%%%%%%%%%%%%%
%             Vertex coloring edge weightings in hypergraphs
%
%                  M K
%                  M K
%                  F P
%
%%%%%%%%%%%%%%%%%%%%%%%%%%%%%%%%%%%%%
\documentclass[reqno,12pt,epsf]{amsart}
\usepackage{amsmath}
\usepackage{latexsym,amsmath}
%\usepackage{mathfont}
%\usepackage[reqno,tbtags]{amsmath}
\usepackage{amsthm}
\usepackage{amssymb}
\usepackage{graphics}

%\textwidth5in
%\textheight7in %\pagestyle{empty} % removes page numbers
%\flushbottom \voffset=-1in \hoffset=-.5in
\newtheorem{theorem}{Theorem}     %[section]
\newtheorem{corollary}[theorem]{Corollary}
\newtheorem{lemma}[theorem]{Lemma}
\newtheorem{fact}[theorem]{Fact}
\newtheorem{proposition}[theorem]{Proposition}
\newtheorem{conjecture}[theorem]{Conjecture}
\newtheorem{remark}[theorem]{Remark}
\newtheorem{definition}[theorem]{Definition}
\newtheorem{question}[theorem]{Question}
\newtheorem{problem}[theorem]{Problem}
\newtheorem{claim}{Claim}




%%%%%%%%%%%%%%%%%%%%%%%%%%%% symbols definitions

\def\rr{\rfloor}  \def\ll{\lfloor}
\def\bi{\bigskip\noindent}
\def\si{\smallskip\noindent}

\def\li{\limits}

\def\de{\delta}

\def\eps{\varepsilon}
\def\om{\omm}
\def\omm{\rho}
%\def\omm{\omega}
\def\al{\alpha}
\def\be{\beta}
\def\R{\mathbb{R}}
\def\N{\mathbb{N}}

\def\per{\mbox{per}}
\def\deg{\mbox{deg}}


\newcounter{rot}
\newcommand{\bignote}[1]{\vspace{0.25 in}\fbox{\parbox{5in}{#1}}%
\marginpar{\fbox{{\bf {\large Q \therot}}}}\addtocounter{rot}{1}}
%%%%%%%%%%%%%%%%%%%%%%%%%%%%%%%%%%%%%%%%%%%%%%%%%%%%%
%%%%%%%%%%%%%%%%%%%%%%%%%%%%%%%%%%%%%%%%%%%%%%%%%%%%%



%\title[Vertex-coloring edge-weightings]
%{Vertex coloring edge weightings in hypergraphs}
\title[The 1-2-3 Conjecture for Hypergraphs]
{The 1-2-3 Conjecture for Hypergraphs}

\author[M. Kalkowski]{Maciej Kalkowski}
\curraddr[M. Kalkowski]{Adam Mickiewicz University\\ Faculty of Mathematics and
 Computer Science\\ Pozna\'n, Poland}
\email{kalkos@atos.wmid.amu.edu.pl}
\author[M. Karo\'nski]{Micha{\l} Karo\'nski}
\curraddr[M. Karo\'nski]{Adam Mickiewicz University\\ Faculty of Mathematics and
 Computer Science\\ Pozna\'n, Poland\\ and Emory University\\
Department of Mathematics and Computer Science\\
Atlanta, GA, USA}
\email{karonski@amu.edu.pl}
\author[F. Pfender]{Florian Pfender}
\curraddr[F. Pfender]{University of Colorado Denver\\ Department of Mathematical and Statistical Sciences\\ Denver, CO, USA}
\email{Florian.Pfender@ucdenver.edu}
\thanks{Research of the third author supported in part by a Simons Foundation Collaboration Grant}
\keywords {irregular hypergraph labelings}
\subjclass {05C78, (05C15)}

\begin{document}


\begin{abstract}
A weighting of the edges of a hypergraph is called vertex-coloring if the weighted degrees of the vertices
yield a proper coloring of the graph, i.e., every edge contains at least two vertices with different weighted degrees. In this paper we show that such a weighting is possible
from the weight set $\{ 1,2 ,\ldots,\max\{5,r+1\}\}$ for all hypergraphs with maximum edge size $r\ge 3$ and not containing edges solely consisting of identical vertices. The number $r+1$ is best possible for this statement.
\end{abstract}

\maketitle

\section{Introduction and Notation}
Regular graphs have been studied in a lot of contexts, and have many properties not shared by other graphs. One may ask what is on the other side of the spectrum, and look for graphs which are as irregular as possible. But what is irregular? It is an easy observation that every graph with at least two vertices contains a pair of vertices of equal degree, so one can not hope for graphs which are {\em totally irregular} in the sense that all vertices have pairwise different degrees. This changes if one considers multigraphs. In fact, by multiplying some edges, one can make every graph totally irregular, as long as the original graph does not contain an isolated edge or two isolated vertices. This observation led to the definition of the {\em irregularity strength} of a graph in~\cite{CJLORS}, the minimum maximum multiplicity one has to use on a given graph.

Later, Karo\'{n}ski, {\L}uczak and Thomason~\cite{KLT} asked a similar question inspired by this concept. What if we do not require {\em all} vertices to have pairwise different degrees, but only require this difference for {\em adjacent} vertices? In other words, we want that the degrees yield a proper vertex coloring. This question led to the so called 1-2-3-Conjecture, stated here in the obviously analogous form using edge weights instead of multiplicities.
\begin{conjecture}\label{123}
For every graph $G$ without isolated edges, there is a weighting $\omm: E(G)\to \{1,2,3\}$,
such that the induced vertex weights $\omm(v):=\sum_{e \ni v}\omm(e)$ properly color $V(G)$.
\end{conjecture}
The 1-2-3-Conjecture is known to be true for several classes of graphs, the best known result for general graphs is by the authors of the current article~\cite{KKP}.
\begin{theorem}\label{125}
For every graph $G$ without isolated edges, there is a weighting $\omm: E(G)\to \{1,2,3,4,5\}$,
such that the induced vertex weights $\omm(v):=\sum_{e \ni v}\omm(e)$ properly color $V(G)$.
\end{theorem}
Shortly thereafter, a total version of the 1-2-3-Conjecture, adaptly called the 1-2-Conjecture, was formulated by Przyby{\l}o and Wo\'{z}niak~\cite{PW}.
\begin{conjecture}\label{12con}
For every graph $G$, there is a weighting $\omm: E(G)\cup V(G)\to \{1,2\}$,
such that the induced total vertex weights $w(v):=\omm(v)+\sum_{e \ni v}\omm(e)$ properly color $V(G)$.
\end{conjecture}
Kalkowski in~\cite{K} came close to settling this conjecture.
\begin{theorem}\label{tot123}
For every graph $G$, there are weightings $\omm: E(G)\to \{1,2,3\}$ and $\omm': V(G)\to \{1,2\}$
such that the induced total vertex weights $w(v):=\omm'(v)+\sum_{e \ni v}\omm(e)$ properly color $V(G)$.
\end{theorem}
One natural and promising approach for both conjectures is the use of Alon's Combinatorial Nullstellensatz (see~\cite{A}). In its most straightforward application, it would prove list versions of the conjectures if successful, leading to the following stronger conjectures, first stated by Bartnicki, Grytczuk and Niwczyk, and by Przyby{\l}o and Wo\'{z}niak, and Wong and Zhu, respectively.
\begin{conjecture}\label{123list}\cite{BGN}
For every graph $G$ without isolated edges, and for every assignment of lists of size $3$ to the edges of $G$, there exists a weighting $\omm: E(G)\to \R$ from the lists,
such that the induced vertex weights $\omm(v):=\sum_{e \ni v}\omm(e)$ properly color $V(G)$.
\end{conjecture}
\begin{conjecture}\label{12list}\cite{PW2},\cite{WZ}
For every graph $G$, and for every assignment of lists of size $2$ to the vertices and edges of $G$, there exists a weighting $\omm: V(G)\cup E(G)\to \R$ from the lists,
such that the induced total vertex weights $w(v):=\omm(v)+\sum_{e \ni v}\omm(e)$ properly color $V(G)$.
\end{conjecture}
Conjecture~\ref{123list} is open even if we allow larger lists of some fixed size $k$. For Conjecture~\ref{12list}, the best result due to Zhu and Wong in~\cite{ZW} generalizes Theorem~\ref{tot123}.
\begin{theorem}\label{tot123list}
For every graph $G$, and for every assignment of lists of size $2$ to the vertices and of size $3$ to the edges of $G$, there exists a weighting $\omm: V(G)\cup E(G)\to \R$ from the lists,
such that the induced total vertex weights $w(v):=\omm(v)+\sum_{e \ni v}\omm(e)$ properly color $V(G)$.
\end{theorem}
All these questions also make sense for hypergraphs when the graph $G$ in the statements is replaced by a hypergraph $H$.
%If we allow edges to contain more than two vertices, we get from graphs to hypergraphs. 
%Both questions above make sense for hypergraphs as well. 
Note that it is easy to construct totally irregular hypergraphs, so the irregularity strength of a hypergraph may actually be $1$ in certain cases. In this manuscript, we want to first consider Conjecture~\ref{123} for hypergraphs. 

To start, we have to decide what we mean by a proper vertex coloring of a hypergraph as there are differing notions. We will consider the weakest notion and call a hypergraph properly colored if it does not contain a monochromatic edge, i.e. an edge containing only vertices from one color class.

Next, we have to classify all hypergraphs which do not allow a vertex coloring edge weighting at all. What is the analogon of an isolated edge in the graph case? We will call a set of vertices of any cardinality {\em twin set} if the vertices in the set are contained in the exact same set of edges. With this notion, it is easy to verify that the only obstacle is an edge consisting of a twin set. In the absence of such edges, a vertex coloring edge weighting with integer weights is always possible. So we will ask for such graphs, what is the minimum maximum edge weight we have to use?

Going from graphs to hypergraphs, one discovers several important classes of hypergraphs invisible in the graph case, we will consider three special classes. A hypergraph is called {\em $k$-uniform} if all its edges have size $k$. If any two edges in a hypergraph intersect in at most one vertex, we call the hypergraph {\em linear} (this property is also called {\em simple} in other places). Note that graphs are exactly the $2$-uniform linear hypergraphs.
A hypergraph is called {\em bipartite} if it allows a proper $2$-coloring. 

In general, we allow multiple edges in our hypergraphs to make some operations easier to state. 
%For two vertices $u,v\in V(H)$, their distance is the minimum $k$ such that there exists a path $u=x_0x_1\ldots x_k=v$, where ${x_i,x_{i+1}}\subseteq e_i$ and $e_i\in E(H)$ for $0\le i\le k-1$.
%
Starting with a hypergraph $H$ with vertex set $V(H)$ and edge set $E(H)$ and a vertex $v\in V(H)$, we define the hypergraph $H-v$ (the deletion of $v$) as the hypergraph with
\begin{align*}
V(H-v)&=V(H)\setminus\{v\},\\
E(H-v)&=\{e\setminus\{v\}:e\in E(H)\}.
\end{align*}
In other words, we delete $v$ from every edge, and we keep the resulting edges.

On the other hand, for $X\subseteq V(H)$, we may consider the {\em induced} hypergraph $H[X]$ with
 \begin{align*}
V(H[X])&=X,\\
E(H[X])&=\{e\in E(H):e\subseteq X\}.
\end{align*}
This time, we only allow edges completely contained in the smaller vertex set.

In the next section we provide some hypergraphs giving lower bounds for a number replacing the $3$ in the 1-2-3-Conjecture. In particular, we show that the statement of the 1-2-3-Conjecture can not be true for general hypergraphs. In fact, it would fail even for linear bipartite hypergraphs.

In the third section, we present the main result of the paper---an upper bound for the hypergraph version of the 1-2-3-Conjecture. We will get a bound for hypergraphs depending linearly on the size $r$ of the largest edge, which matches our lower bound as long as $r\ge 4$. %For non-linear hypergraphs our bound is weaker, but still linear.
In Section~\ref{total}, we will turn to Conjecture~\ref{12con}, and extend Theorem~\ref{tot123} to hypergraphs.

%In a final section, we consider the related topic of total weightings.

\section{Lower Bounds}
Let $F$ be any hypergraph with vertex set $V(F)$ and edge set $E(F)$, and suppose that $F$ has minimum degree at least $2$ and no edges of size smaller than $2$. From this, we create another hypergraph $H$, the dual of the incidence graph of $F$, as follows. Let $V(H)$ consist of the vertex-edge incidences in $F$, i.e., pairs $(v,e)$ where $v\in V(F)$, $e\in E(F)$ and $v\in e$. Let 
\begin{align*}
E_1(H)&= \bigcup_{v\in V(F)}\{\{(v,e)\in V(H):e\in E(F)\}\},\\
E_2(H)&= \bigcup_{e\in E(F)}\{\{(v,e)\in V(H):v\in V(F)\}\},\\
E(H)&=E_1(H)\cup E_2(H).
\end{align*}
With this construction, $H$ is linear and $2$-regular, the largest edge in $E_1(H)$ has size equal to the maximum degree in $F$, and the largest edge in $E_2(H)$ has size equal to the largest edge in $F$. 
%Note that for every shortest path $(v_0,e_0)(v_1,e_1)\ldots (v_k,e_k)$ in $H$ and $0\le i<j\le k$, we have 
%\[|\{v_i,e_i\}\cap\{v_{j},e_{j}\}|=
%\begin{cases}
%1,&\mbox{if }j=i+1,\\
%0,&\mbox{if }j>i+1.
%\end{cases}
%\]
%
%As $H$ does not contain any odd cycles, $H$ is bipartite: start with any vertex $(v,e)\in V(H)$, and partition the remaining vertices
%in the same component by the parity of their distance to $(v,e)$. 
Further, $H$ is bipartite as we can properly $2$-color $H$ as follows. Let $G$ be the graph underlying $H$, i.e. $V(G)=V(H)$ and $(v,e)(v',e')\in E(G)$ if and only if $|\{v,e\}\cap\{v',e'\}|=1$. As $H$ is linear, we can greedily find a subgraph $G'\subseteq G$ such that every edge in $E(H)$ contains exactly one edge in $G'$. As $H$ is $2$-regular, $G'$ has maximum degree at most $2$. Further, the edges of $G'$ naturally partition into two matchings $E_1(G')$ and $E_2(G')$ by their association with $E_1(H)$ and $E_2(H)$. Therefore, $G'$ has no odd cycles, and we can find a proper 2-vertex-coloring of $G'$, which is also a proper 2-vertex-coloring of $H$ by construction.

Further, suppose that $F$ has chromatic number $\chi(F)$, and consider any weighting $\omm: E(H)\to \{1,2,\ldots,\chi(F)-1\}$.  
Note that every edge $h_1=\{(v,e_1),(v,e_2),\ldots,(v,e_r)\}\in E_1(H)$ corresponds to the vertex $v\in V(F)$, and every edge 
$h_2=\{(v_1,e),(v_2,e),\ldots,(v_s,e)\}\in E_2(H)$ corresponds to the edge $e\in E(F)$. With this correspondence, we can define 
$\omm(v):=\omm(h_1)$ and $\omm(e):=\omm(h_2)$, and then
\[
\omm((v,e))=\sum_{h\ni (v,e)}\omm(h)=\omm(v)+\omm(e).
\]
Thus, the edge $h_2$ from above is monochromatic if and only if all the $\omm(v_i)$ are the same. If $\omm$ induces a proper coloring on the vertices of $H$, then $\omm$ has to be a proper coloring on the vertices of $F$, a contradiction. Therefore, $\omm$ cannot induce a proper vertex coloring of $H$.

This construction gives us several lower bounds. If we start with a complete graph $F$ on $r+1$ vertices, we obtain a hypergraph with maximum edge size $r$,  which needs a weight set of at least $\{1,2,\ldots,r+1\}$ on the edges to properly color the vertices. Starting with any other $r$-regular graph $F$ with chromatic number $r$,  we obtain a hypergraph with maximum edge size $r$,  which needs a weight set of at least $\{1,2,\ldots,r\}$ on the edges to properly color the vertices.

If we start with the Fano plane (or any other $3$-regular $3$-uniform non-bipartite hypergraph), we obtain a $2$-regular $3$-uniform hypergraph, which needs a weight set of at least $\{1,2,3\}$ on the edges to properly color the vertices. 

On the other hand, this construction cannot give us non-trivial examples for $r$-uniform hypergraphs with $r\ge 4$.  Thomassen shows in~\cite{T} that all $r$-uniform $r$-regular hypergraphs are bipartite for $r\ge 4$, leaving open the possibility of a vertex coloring edge weighting from the set $\{1,2\}$ for such graphs.

\section{Upper Bounds}

For $r\ge 4$, we show that the bound of $r+1$ we provided in the last section is in fact best possible. %Notice that for linear hypergraphs without multiple edges, edges consisting of a twin set are exactly isolated edges and edges of size at most $1$.
\begin{theorem}\label{Thm:Main}
For every hypergraph $H$ with edges of order at most $r\ge 2$, and no edge consisting of a twin set, there is a weighting $\omm: E(H)\to \{1,2,\ldots,\max\{5,r+1\}\}$,
such that the induced vertex weights $\omm(v):=\sum_{e\ni v}\omm(e)$ properly color $V(H)$.
\end{theorem}

\begin{proof}
We prove the statement by induction on $n=|V(H)|$. 
In fact, we will prove a slightly stronger statement to make the induction work. The statement is stronger because we can pick $\omm'$ to be constant:\\

{\em 
For every hypergraph $H$ with all edges of order between $2$ and $r\ge 2$, and no edge consisting of a twin set, and for every weighting of the vertices $\omm ':V(H)\to \N$, there is a weighting $\omm: E(H)\to \{1,2,\ldots,\max\{5,r+1\}\}$,
such that the induced vertex weights $\omm(v):=\omm '(v)+\sum_{e\ni v}\omm(e)$ properly color $V(H)$.
}\\

The statement is easy for $n=3$, so assume that $n\ge 4$.
We may assume that there are no twin sets of size larger than 1, otherwise we can remove all but one of the vertices of such a twin set and use induction. Further, we may assume that there are no multiple edges, as we can fix a weight on one of them and forget about the edge in the further argument.
We may also assume that every vertex lies in an edge of size $2$. Otherwise, pick a vertex $v$ which is in no edge of size $2$, and consider $H-v$. The only edges changed by this are edges containing $v$. If one of these edges now consists of a twin set, then that twin set was a twin set of size larger than $1$ in $H$, a contradiction.
Any edge weighting inducing a proper coloring on $H-v$ then induces a proper coloring on $H$ as well. 
%Note that this also implies that $H$ contains no pair of twins, as these would have to be in a common $2$-edge, and by linearity this edge would be isolated.
%Further, we may assume that H$ is connected, otherwise we can argue component-wise.

The main idea of the proof is as follows. We order the vertices in a specific ordering and then only consider an associated graph on the same vertex set, and  all edges consisting of the first two vertices in each edge of $H$. Then we proceed very similarly to the proof in~\cite{KKP} to weight these edges respecting the order of the vertices, guaranteeing that in the end, the first two vertices of each edge in $H$ have different weighted degrees. Minor shortcomings we can fix in the end.
The main difficulty in this approach lies in the fact that in the proof for graphs we must pick a vertex ordering with specific properties, but as the associated graph of the hypergraph depends on the ordering of the vertices, we cannot change this vertex ordering after restricting our attention to the associated graph. To circumvent this problem,
we very carefully pick the ordering such that the resulting associated graph already has properties very close to what we need to make the graph process work without reordering.


In this spirit, define $E_2$ to be the set of edges of size $2$.
Let $E_1\subseteq E(H)$ be the set of edges in $H$ which contain an edge of $E_2$, and let $E_0=E(H)\setminus (E_1\cup E_2)$. %Let $d_1(v)$ ($d_2(v)$) stand for the number of edges in $E_1(H_n)$ ($E_2(H_n)$) incident to $v$.
For any ordering $\pi$ of the vertices, let $E_\pi$ be the set of pairs of vertices appearing first and second with respect to $\pi$ in an edge of size at least $3$ in $H$. Now let us find a suitable ordering $\pi$, building backwards from the last vertex in the ordering.

If $H$ contains a vertex incident to at least two edges in $E_2$, make such a vertex the last vertex $v_n$, 
and pick $v_{n-1},v_{n-2}$ such that $v_{n-1}v,v_{n-2}v\in E_2$. Let $h=v_{n-2}v$, $t=2$,
and skip forward to the next paragraph of this construction. 
Now assume that $H$ contains no such vertex, so $E_2$ is a perfect matching. 
If $E_0=\emptyset$, then the proper weighting is easy. Changing the weight of an edge in $E_1$ does not impact the properness of the coloring of the contained edge in $E_2$. Since $H$ has no edge consisting of a twin set, every edge in $E_2$ has an intersection of exactly one vertex with at least one edge in $E_1$.
%, and we can greedily make all edges in $E_2$ properly colored by changing the weights of the edges in $E_1$ only. 
Note that every edge in $E_1$ can intersect at most $r-1$ edges in $E_2$, so we can greedily pick a weight for each edge in $E_1$ such that all edges in $E_2$ intersecting that edge in exactly one vertex are properly colored.
Since all edges in $E_1$ contain edges in $E_2$, this in turn makes the coloring of every edge in $E_1$ proper.
Thus, we may assume that $E_0\ne \emptyset$. 
Find a hyperedge $h\in E_0$ of minimal size $3\le t\le \frac{n}2$, and choose $\pi$ such that $h=\{v_n,v_{n-2},\ldots,v_{n-2t+2}\}$, and $v_{n-2i}v_{n-2i-1}\in E_2$ for $0\le i\le t-1$.


Then, successively for $i\ge 2t$, let $v_{n-i}$ be a vertex not in the set of previously selected vertices $\{v_{n-i+1},\ldots,v_n\}$ connected to $\{v_{n-i+1},\ldots,v_n\}$ by an edge in $E_2\cup E_\pi$, as long as  such a vertex exists. Note that at this point, $\pi$ is determined sufficiently to decide if there is an edge in $E_\pi$ between $\{v_{n-i+1},\ldots,v_n\}$ and the remaining vertices. If we arrive at $v_1$ this way, $\pi$ is determined. If the process stops before, say after assigning $i$ labels, delete the previously ordered vertices from $H$ to form a hypergraph $H'=H-\{v_{n-i+1},\ldots,v_n\}$ on $n'=n-i$ vertices. Clearly, $H'$ contains no edges of size $1$. As $H$ contains no twin set of size $2$, $H'$ contains no twin set of size $2$ either, and therefore no edges consisting of twin sets. By induction, we can find a vertex coloring edge weighting on $H'$. Similarly, let $H''=H[\{v_{n-i+1},\ldots,v_n\}]$ be the connected hypergraph induced by $H$ on the previously selected vertices. Add the weights of edges intersecting both $V(H')$ and $V(H'')$ which we computed in the weighting of $H'$ to the respective vertex weights in $H''$, and use induction to weight the edges in $H''$, finishing the proof. Thus, we may assume in the following that $\pi$ is completely determined, and let 
 $G=G_\pi$  be the graph with edge set $E_2\cup E_\pi$.

%Let $G_1=G[\{v_1,v_2\ldots,v_{n-2t+4}\}]$, and notice that $G_1$ is connected.
Notice that $G$  is ordered in a way that every vertex up to $v_{n-2t+3}$ has a neighbor later in the order.
 
 Now we assign weights very similarly to the proof in~\cite{KKP}, we repeat large parts of this proof here so that this article is self contained. When we assign a weight to an edge in $G$, we are assigning it at the same time to the edge in $H$ that corresponds to the edge in $G$. 

%==================

%We may assume that $G$ is connected, otherwise we can argue component-wise. If $|G|\le 1$, there is nothing to prove, so we may assume that $|G|\ge 3$ and $G$ contains a vertex with degree $\ge 2$. Order the vertices $V(G)=\{ v_1,v_2,\ldots,v_n\}$ so that $d(v_n)\ge 2$ and for every $1\le i\le n-1$, $v_i$ has a neighbor in $\{ v_{i+1},v_{i+2},\ldots, v_n\}$.

%Let $G[\{v_1,\ldots,v_s\}]$ be the first component of $G$, where $s=n$ or $s=n-2t+4$.
%Let $v_s$ be the last vertex in $G_1$, i.e., $s\in \{n-2t+4,n\}$. 
We start by assigning the provisional weight $\om(e)=3$ to every edge $e$ in $G$ and adjust it at most twice while going through the vertices in the order specified by $\pi$---once when we are considering the first vertex in the edge, and once when we consider the second vertex. To every vertex $v_i$ with $i<n-2t+2$, we will assign a set of two colors $W(v_i)=\{w(v_i), w(v_i)+2\}$ with $w(v_i)\in \{ 0,1\}\mod 4$, so that for every edge $v_jv_i\in E(G)$ with $1\le j<i$, we have $W(v_j)\cap W(v_i)=\emptyset$, and we will guarantee that $\om(v_i)=\om'(v_i)+\sum_{e\ni v_i}\om(e)\in W(v_i)$. 
%Finally we will adjust the weights of the edges incident to $v_s$ in $G$ to make sure that $\om(v_s)$ is different from $\om(v_i)$ for all $v_i\in N_{G_1}(v_s)$.

To this end, let $\om(v_1)=\omm '(v_1)+3d_H(v_1)$, and pick the set $W(v_1)=\{ w(v_1),w(v_1)+2\}$ so that $\om(v_1)\in W(v_1)$ and $w(v_1)\in \{ 0,1\}\mod 4$.
Let $2\le k\le n-2t+1$ and assume that we have picked $W(v_i)$ for all $i<k$ and
\begin{itemize}
\item $\om(v_i)\in W(v_i)$ for $i<k$,
\item $\om(v_kv_j)=3$ for all edges in $G$ with $j>k$, and
\item if $\om(v_iv_k)\ne 3$ for some edge in $G$ with $i<k$, then $\om(v_iv_k)=2$ and $\om(v_i)=w(v_i)$ or $\om(v_iv_k)=4$ and $\om(v_i)=w(v_i)+2$.
\end{itemize}
%
% (and $\om(v_i)\in W(v_i)$), all edges $v_kv_j$ with $k<j$ still have $\om(v_kv_j)=3$, and if $\om(v_jv_k)\ne 3$ for $j<k$, then $\om(v_jv_k)=2$ and $\om(v_j)=w(v_j)$ or $\om(v_jv_k)=4$ and $\om(v_j)=w(v_j)+2$.

If $v_iv_k\in E(G)$ for some $i<k$ we can either add or subtract $2$ to $\om(v_iv_k)$ keeping $\om(v_i)\in W(v_i)$. If $v_k$ has $d$ such neighbors, this gives us a total of $d+1$ choices (all of the same parity) for $\om(v_k)$. In addition to this we will allow to alter the weight $\om(v_kv_j)$ by $1$, where $j>k$ is smallest such that $v_kv_j\in E(G)$. This way, $\om(v_k)$ can take all integer values of an interval $[a,a+2d+2]$. We want to adjust the weights and assign $w(v_k)$ so that
\begin{enumerate}
\item $\om(v_i)\in W(v_i)$ for $1\le i\le k$,
\item $w(v_i)\ne w(v_k)$ for $v_iv_k\in E(G)$ with $i<k$, and\label{con1}
\item either $\om(v_k)=w(v_k)$ and $\om(v_kv_j)\in \{2,3\}$ or $\om(v_k)=w(v_k)+2$ and $\om(v_kv_j)\in \{3,4\}$.\label{con2}
\end{enumerate}
Condition~(\ref{con1}) can block at most $2d$ values in $[a,a+2d+2]$, and condition~(\ref{con2}) can block only the values $a$ and $a+2d+2$ (for all other values $\om(v_k)$ with $\om(v_kv_j)\ne 3$, we have the choice between $\om(v_kv_j)=2$ and $\om(v_kv_j)=4$). At least one value remains open for $\om(v_k)$.

This way, we can assign the sets $W(v_k)$ step by step for all $k\le n-2t+1$ without conflicts. Note that the first time $\om(v_k)$ may get changed by an adjustment of an edge $v_kv_i$ for $i>k$ is when $i=j$, so we don't run into problems with edges weighted $2$ or $4$.

%As the final step, we have to find an open weight for $v_s$. 
%%This finalizes $\om$ and finishes the proof by setting $\omm=\om$.
%
%If $s=n$, then this finalizes the weighting of the edges, and we have only used edge weights from the set $\{1,2,\ldots,5\}$.
%%, and we have not used the fact that $H$ is linear in this part of the proof.

As the final step, we make sure that the remaining edges are also colored properly by the weighting. 
We consider the cases $t=2$ and $t\ge 3$ separately, starting with $t=2$.
In this case, continue the previous process until $v_{n-1}$ in $G$, which is possible as both $v_{n-2}$ and $v_{n-1}$ are adjacent to $v_n$. Now we have to find an open weight for $v_n$.
%This time, we don't have an extra edge to work with, but we don't have to worry about later vertices.
If $v_iv_n\in E(G)$ for some $i<n$ we can again either add or subtract $2$ to $\om(v_iv_n)$ keeping $\om(v_i)\in W(v_i)$.
% Note that, for $i < n$, if $f(v_i)\in{0,1} \mod 4$ then we can add $2$ to $f(v_iv_n)$ and if $f(v_i)\in {2,3} \mod 4$ then we can subtract $2$. 
These possible adjustments give a total of $d_{G}(v_n)+1\ge 3$ options (all of the same parity) for $\om(v_n)$. Hence if the smallest such
% Adjusting the edges incident to $v_n$ by $2$ like above gives us a total of $d(v_n)+1\ge 3$ options (all of the same parity) for $\om(v_n)$.
% If the smallest such
option $a$ has $a\in\{2,3\}\mod 4$, then picking the lower possible weight on each edge incident to $v_n$ gives a proper coloring of the vertices. If $a\in\{0,1\}\mod 4$ and there is a $v_i\in N(v_n)$ with $w(v_i)\ne a$, then picking the higher weight on $v_iv_n$ and the lower weight on all other edges gives $\om(v_n)=a+2$ in a proper coloring. Finally, if $a\in\{0,1\}\mod 4$ and $w(v_i)= a$ for all $v_i\in N(v_n)$, picking the higher weight on at least two edges gives a proper coloring.



%In fact, we will stop our weighting procedure a little bit earlier, after we have considered $v_{s-3}=v_{n-2t+1}$, the predecessor of the first vertex of $h$.
For $t\ge 3$, observe that after dealing with $v_{n-2t+1}$, every edge in $H$ which contains at least two vertices in $\{v_1,v_2,\ldots,v_{n-2t+1}\}$ is properly colored, and we will not change the weights of these vertices in the remainder of the proof. Every remaining edge in $E_1$ contains at least one of the pairs $\{v_{n-2i},v_{n-2i-1}\}$ for $0\le i\le t-1$. Every remaining edge in $E_0$ contains at least $t-1$ vertices in $\{v_{n-2t+2},\ldots,v_n\}$ by the minimality of $t=|h|$, and therefore at least one of the pairs $\{v_{n-2t+2},v_{n-3}\}$, $\{v_{n-2t+2},v_{n-2}\}$,  $\{v_{n-2t+2},v_{n-1}\}$, $\{v_{n-2t+2},v_{n}\}$, $\{v_{n-3},v_{n-1}\}$, $\{v_{n-3},v_{n}\}$, $\{v_{n-2},v_{n-1}\}$,  $\{v_{n-2},v_{n}\}$. Therefore, if we adjust the weights of the edges $h$, $v_{n-3}v_{n-2}$ and $v_{n-1}v_n$ such that all these $t+8$ pairs receive different weights, our weighting is complete.

First, adjust $h$ such that all of the pairs $\{v_{n-2i},v_{n-2i-1}\}$ for $0\le i\le t-1$ are colored properly. As we have at least $r+1\ge t+1$ weights to choose from, this is possible. 
From each of the at least five choices for  $\om(v_{n-3}v_{n-2})$ and $\om(v_{n-1}v_{n})$, at least three choices will guarantee that $\om(v_{n-2t+2})\ne \om(v_{n-i})$ for $0\le i\le 3$. So let $x_1<x_2<x_3$ be such remaining choices for $\om(v_{n-3}v_{n-2})$, and let $y_1<y_2<y_3$ be such remaining choices for $\om(v_{n-1}v_{n})$.
Notice that $y_1-x_3<y_1-x_2<y_1-x_3<y_2-x_3<y_3-x_3$, so we have at least five different options for $\om(v_{n-1}v_{n})-\om(v_{n-3}v_{n-2})$ at this point. Therefore, at least one of these options will yield a result with $\om(v_{n-i})\ne \om(v_{n-j})$ for all $0\le i\le 1$, $2\le j\le 3$. This proves the theorem.
\end{proof}



% Since isolated edges in $G_\pi$ which may appear cause more problems if the graph is not linear, we will consider these edges first in the weighting, and later make sure that these edges do not become monochromatic. 
%
%We start by finding a similar ordering $\pi$ as in the proof of Theorem~\ref{Thm:Main}. In the case that two edges in $E_2$ intersect, we use the exact same ordering until there are no edges in $G_{\pi}$ between vertices already put into the ordering and the remaining vertices. If we do not reach $v_1$ this way on the first try, we can again use induction to show the existence of the weighting in the theorem statement. If we do reach $v_1$, we will find a weighting $\omm: E(H)\to \{1,2,3,4,5\}$ as stated in the proof of Theorem~\ref{Thm:Main}.
%
%In the case that $E_2$ is a matching, we cannot guarantee the existence of a hyperedge $h=\{v_{n-2t+2},v_{n-2t+4},\ldots,v_n\}$. Instead, we will use a hyperedge $h$ which intersects some edge in $E_2$ in exactly one vertex, and, given that constraint, intersects the minimal number of edges in $E_2$ (say a total of $t$ intersected edges). Now, as above, put the edges in $E_2$ intersecting $h$ last in the ordering $v_{n-1}v_n,v_{n-3}v_{n-2},\ldots, v_{n-2t+1}v_{n-2t+2}$ with $\{v_{n-2t+2},v_{n-2t+4},\ldots,v_n\}\subseteq h$ and $v_1,v_2,\ldots,v_{n-2t+1}\notin h$. Again, extend this ordering backwards as far as possible, guaranteeing that every vertex added to the order has a neighbor in $G_\pi$ later in the ordering.
%
%If this process stops before we reach $v_1$, we cannot use the same inductive argument as before---the reasons will become clear very soon. Instead, we delete all previously ordered vertices and continue with the remaining hypergraph until all vertices are in order. Note that the remaining hypergraph does not contain edges of size 1.
%With this order, every component of $G$ consists of a set of consecutive vertices in $\pi$. Note that the component containing $v_1$ contains more than $2$ vertices as otherwise $v_1$ would be isolated or $v_1$ and $v_2$ would be twins.
%%, resulting in more than one component with more than two vertices, and possibly several components with exactly two vertices.
%%Again, we may assume by induction that we find $\pi$ on the first try. Further, we may assume that $E_2$ is a matching.
%
%Now give every edge a preliminary weight of $2r-1$. Next, we change the weight of some edges in a way that all edges in $E_2$ intersecting $h$ other than possibly $v_{n-2t+1}v_{n-2t+2}$ have two weights on their vertices which are different modulo $r-1$. In the final step we will only adjust weights of edges by multiples of $r-1$, and in this way these preprocessed edges can never become monochromatic.
%
%For this, let $E^*$ contain all edges in $E_2$ consisting of the last two vertices of a component of $G$ in the ordering given by $\pi$, unless that component ends with a vertex with at least two different neighbors in $G$. Note that $E^*$ includes all edges in components of size $2$.
%As none of these edges consists of a twin set of vertices, for every $f\in E^*$ there exists an edge $e_f\in E(H)$ containing exactly one of the two vertices of $f$. Every edge $e\in E(H)$ can have a unique intersection with at most $r-2$ edges $f\in E^*$: If $e$ intersects $r$ or $r-1$ different edges in $E_2$ uniquely, then $e$ does not fully contain an edge in $E_2$, and so the first two vertices of $e$ are not in $E^*$. Thus, we can greedily add the lowest suitable values from $\{0,\ldots,r-2\}$ to the weights of the edges $e_f$ one-by-one, making the two vertex weights in each edge $f\in E^*$ different modulo $r-1$.
%
%Again, we proceed with the same weighting algorithm as in the proof of Theorem~\ref{Thm:Main}. But instead of adjusting edge weights by $\{-2,-1,0,1,2\}$, we adjust them by $\{-2(r-1),-(r-1),0,(r-1),2(r-1)\}$ this time. Other than this extra factor of $(r-1)$, the weighting algorithm works exactly the same. Since we preprocessed  the edges in $E^*$, these adjustments will never make the edges in $E^*$ monochromatic. 
%Again, by induction it is enough to look at one component of size at least $3$ of $G$ now, say the last. With the same argument as in the proof of Theorem~\ref{Thm:Main}, we can find the weighting if $E_2$ is not a matching. So we assume that $E_2$ is a matching, and the last components of $G$ are edges of this matching. The edge $h\in E(H)$ defined earlier in this proof intersects the last $t$ edges of this matching.
%
% We stop the algorithm after adjusting the edges around $v_{n-2t+2}$. Note that the biggest possible change required on $h$ by the algorithm at this step is $(r-1)$. Further, the weight of $h$ could have been adjusted upward by up to $(r-2)$ in the preprocessing step dealing with $E^*$. So $h$ has weight between $(2r-1)-(r-1)=r$ and $(2r-1)+(r-2)+(r-1)=4r-4$ at this point. The edges in $E^*$ are properly colored due to the preprocessing, all other edges other than $h$ are properly colored due to the algorithm. To make $h$ properly colored, we can adjust the weight of $v_{n-1}v_n$ by $(r-1)$, which does not affect the properness of any other edge at this point.
%
%After the algorithm runs, every edge has weight at least $(2r-1)-2(r-1)=1$ and at most $(2r-1)+(r-2)+2(r-1)=5r-5$, proving the theorem. 
%\end{proof}

\section{Total Colorings}\label{total}
In this short section, we generalize Theorem~\ref{tot123}  % and~\ref{tot123list} 
to hypergraphs. While a generalization of the stronger Theorem~\ref{tot123list} to hypergraphs would be very interesting, we are not able to prove such a statement at this point.

\begin{theorem}
For every hypergraph $H$ without edges of size $0$ or $1$, there are weightings $\omm: E(H)\to \{1,2,3\}$ and $\omm': V(H)\to \{1,2\}$
such that the induced total vertex weights $w(v):=\omm'(v)+\sum_{e\ni v}\omm(e)$ properly color $V(H)$.
\end{theorem}

\begin{proof}
In comparison with the previous theorem, the proof of Theorem~\ref{tot123} generalizes to hypergraphs with even fewer complications. To this end, order the vertices of $H$ in any order $\pi$, and consider the multigraph $G=G_{\pi}$ as above, i.e. the graph on the same vertex set with edges $u_ev_e$, where $u_e$ and $v_e$ are the first and second vertex appearing in some edge $e\in E(H)$. We will show the stronger statement that we can find a weighting such that $w(u_e)\ne w(v_e)$ for every edge $e\in E(H)$.

%Now proceed with the proof of Theorem~\ref{tot123} on $G$, taking the whole edges of $H$ into account. %, i.e.:
%\begin{enumerate}
%\item 
Start with all edges weighted $2$, and all vertices weighted $1$ to give a temporary total weight to every vertex.
%\item 
Change the weights of edges and vertices greedily going through the vertices in the order given by $\pi$. When we consider a vertex $v$, the only edge weights we may adjust are those of edges $e$ with $v_e=v$, and the vertex weights of $u_e$ and $v_e$ in these edges.
%\item 
Note that we can adjust the weight of every such edge $e$ by either $+1$ or $-1$, and change the weight of $u_e$ at the same time to keep the total weight of $u_e$ unchanged, always using only allowable weights.
%\item 
This gives us sufficiently many options for the total weight of $v$ to make sure that this total weight is different from all the total weights of  the vertices $u_e$ with $v_e=v$ (if there are $d$ such vertices, we have $d+1$ different options for the total weight of $v$).
As the total weight of $v$ is fixed for the remainder of the process, we end up with a weighting with the desired properties.
%\end{enumerate}
\end{proof}


%We will not present the proof here, as the statement is implied by the list version.
%\begin{theorem}\label{main}
%For every hypergraph $H$ without edges of size $0$ and $1$, and for every assignment of lists of size $2$ to the vertices and of size $3$ to the edges of $G$, there exists a weighting $\omm: V(G)\cup E(G)\to \R$ from the lists,
%such that the induced total vertex weights $w(v):=\omm(v)+\sum_{e\ni v}\omm(e)$ properly color $V(G)$.
%\end{theorem}
%\begin{proof}
%The proof is an adaptation of the proof of Theorem~\ref{tot123list} to hypergraphs, using some ideas from our previous proofs.
%Order the vertices in some order.
%For every edge $e\in E(H)$, let $u_e,v_e\in e$ be the first two vertices in the edge. We will show the stronger statement, that we can find a weighting such that $w(u_e)\ne w(v_e)$ for every edge $e\in E(H)$.
%We may assume that $\{u_e,v_e\}\ne \{u_{e'},v_{e'}\}$ for any pair of edges $e,e'\in E(H)$. Otherwise, pick a weight from the list of $e'$, and add it to the lists of all vertices in $e'$. If we can now find a coloring weighting in the Hypergraph $H-e'$ from the new lists, the corresponding weighting in $H$ will also induce a proper coloring.
%
%Now consider the polynomial
%\[
%\phi(\om)=\prod_{e\in E(H)}(w(u_e)-w(v_e)).
%\]
%Note that non-zeros of $\phi$ correspond exactly to total weightings of $H$ with $w(u_e)\ne w(v_e)$ for every edge $e\in E(H)$. Let $n=|V(H)|$, $m=|E(H)|$, and let $A_H$ be an $m\times (m+n)$ matrix, where the rows are labeled with the elements of $E(H)$, the columns are labeled with the elements of $E(H)\cup V(H)$, and
%\[
%A_H(e,z)=
%\begin{cases}
%~1,&\mbox{ if }z\in V(H)\mbox{ and }z\in e,\\
%~1, &\mbox{ if }z\in E(H)\mbox{ and }u_e\in z,v_e\notin z,\\
%-1, &\mbox{ if }z\in E(H)\mbox{ and }v_e\in z,u_e\notin z,\\
%~0, &\mbox{ otherwise}.
%\end{cases}
%\]
%Let $B$ be an $m\times m$ matrix consisting of columns of $A_H$, such that every column corresponding to a vertex of $H$ appears at most once, and every column corresponding to an edge of $H$ appears at most twice. Then the permanent $\per(B)$ equals a coefficient of a maximum degree monomial in $\phi$. By a standard application of Alon's Combinatorial Nullstellensatz, there exists a non-zero of $\phi$ (and thus a proper coloring by total vertex weights from the lists) from any given list assignment with vertex lists of size two and edge lists of size three, if we can find such a matrix $B$ with $\per(B)\ne 0$.
%
%We will find such a $B$ by induction on $n$. For $n=0$, the statement is trivial as by definition, the permanent of a $0\times 0$ matrix is $1$. For $n\ge 1$, let $u\in V(H)$ be the first vertex in the order, and consider the hypergraph $H'$ induced by $H$ on $V(H)\setminus\{u\}$, i.e., $H'$ contains all edges of $H$ which do not contain $u$. Let $k=\deg(u)$, then by induction there is an $(m-k)\times (m-k)$-matrix $B'$ with $\per(B')\ne 0$ consisting of columns of $A_{H'}$, at most two columns equal to $A_{H'}(.,e)$ for each edge, and at most one column equal to $A_{H'}(.,v)$ for each vertex.
%
%Now build an $m\times m$ matrix $C$ from $B'$ by the use of the corresponding columns of $A_H$, and the addition of $k$ identical columns equal to $A_H(.,u)$. These added columns all have $k$ entries equal to $1$, and the remaining entries $0$, so $\per(C)=k!~\per(B')\ne 0$. Now, whenever $u\in e\in E(H)$, and $C$ contains $A_H(.,v_e)$, replace that column by the column $A_H(.,e)$. Note that $A_H(.,e)=A_H(.,u)-A_H(.,v_e)$. By the multilinearity of the permanent, the difference of permanents of $C$ and the matrix after the switch is equal to the permanent of a matrix containing $(k+1)$ copies of $A_H(.,u)$, a singular matrix with permanent $0$. Thus, we can make all these switches one-by-one, arriving at a matrix $D$ without columns equal to $A_H(.,v_e)$ and at most one column equal to $A_H(.,e)$ for $u\in e\in E(H)$, and with $\per(D)=\per(C)$.
%
%To now get $B$ from $D$, we will replace the columns equal to $A_H(.,u)$ appropriately one-by-one. For every $e\ni u$, replacing one column $A_H(.,u)$ by one of $A_H(.,v_e)$ and $A_H(.,e)$ will result in a matrix with $\per\ne 0$. If both the resulting matrices had permanent $0$, then by the multilinearity of the permanent, the matrix before the replacement would have had permanent $0$, a contradiction. After replacing all columns equal to $A_H(.,u)$ this way, we arrive at $B$ with the desired properties.
%\end{proof}

\section{Conclusion and Open Questions}

The only class of hypergraphs we know achieving the bound in Theorem~\ref{Thm:Main} is the one constructed above, stemming from the complete graphs $K_{r+1}$. Possibly, it is true that this is the unique example for $r\ge 3$, and in all other cases a set $\{1,2,\ldots,r\}$ is sufficient.

Note that most of our examples on the lower bounds are highly non-uniform, they contain very small and very large edges. 
For $r$-uniform hypergraphs, there may be a constant upper bound instead, independent of $r$. For $r=3$, we conjecture the following.
\begin{conjecture}\label{con3}
For every $3$-uniform hypergraph $H$ without an isolated edge, there is a weighting $\omm: E(H)\to \{1,2,3\}$,
such that the induced vertex weights $\omm(v):=\sum_{e\ni v}\omm(e)$ properly color $V(H)$.
\end{conjecture}
At some point we thought that it may even be true that for $r\ge 4$, the set $\{1,2\}$ is sufficient, and in fact asked this question in a previous version of this manuscript. 
Then Bennett, Dudek, Frieze and Helenius~\cite{BenDFH} constructed a family of  uniform hypergraphs for all $r$ which require the set $\{1,2,3\}$, and extended Conjecture~\ref{con3} accordingly. 
\begin{conjecture}\cite{BenDFH}
For every $r$-uniform hypergraph $H$ without an isolated edge, there is a weighting $\omm: E(H)\to \{1,2,3\}$,
such that the induced vertex weights $\omm(v):=\sum_{e\ni v}\omm(e)$ properly color $V(H)$.
\end{conjecture}
The construction in~\cite{BenDFH} is as follows. Start with a cycle of length $4k+2$, and transform every $2$-edge into an $r$-edge by producing $(r-2)$ twins of every other vertex in the cycle. As the cycle cannot be properly colored by an edge weighting from $\{1,2\}$, this is impossible for this hypergraph as well. If we forbid this type of construction by forbidding twins, we think the following is true.
\begin{conjecture}
For large enough $r$, for every $r$-uniform hypergraph $H$ without twins, there is a weighting $\omm: E(H)\to \{1,2\}$,
such that the induced vertex weights $\omm(v):=\sum_{e\ni v}\omm(e)$ properly color $V(H)$.
\end{conjecture}

%\begin{conjecture}
%For $r$ large enough, for every $r$-uniform hypergraph $H$ without an isolated edge, there is a weighting $\omm: E(H)\to \{1,2\}$,
%such that the induced vertex weights $\omm(v):=\sum_{e\ni v}\omm(e)$ properly color $V(H)$.
%\end{conjecture}

%As mentioned above, list versions of all the above questions may be asked

%Another interesting question is the list version of the problem. Instead of choosing from weights $\{ 1,2,\ldots,k\}$ on each edge, each edge has its own list of $k$ different available weights. A natural strategy to attack the 1-2-3-Conjecture is Alon's Combinatorial Nullstellensatz (see~\cite{A}), and in its most basic application one would prove the list version. For graphs, this approach has been successful only for special graph classes, for general graphs no sufficient $k$ is known.
%
%One area where the Nullstellensatz has been used successfully is a related problem on total weightings. Here, the approach used for graphs by Zhu and Wong~\cite{ZW} can also be carried over to hypergraphs, see~\cite{P}.

 \bibliographystyle{amsplain}
 \begin{thebibliography}{99}
%\bibitem

\bibitem{A} N. Alon, Combinatorial Nullstellensatz, { Combin. Prob. Comput.} { 8} (1999), 7--29.  
\bibitem{BGN} T. Bartnicki, J. Grytczuk, S. Niwczyk, Weight choosability of graphs, { J. Graph Theory} { 60}, (2009), 242--256.
\bibitem{BenDFH} P. Bennett, A. Dudek, A. Frieze, L. Helenius, The 1-2-3 Conjecture for Uniform Hypergraphs, preprint, arXiv:1511.04569v1.
\bibitem{CJLORS} G. Chartrand, M.S. Jacobson, J. Lehel, O.R.
Oellermann, S. Ruiz, and F. Saba, Irregular Networks, { Congressus
Numerantium} { 64} (1988), 187--192.
%\bibitem{D} R. Diestel, ``{\em Graph Theory},'' Springer Verlag, Heidelberg, 2005.
\bibitem{K} M. Kalkowski, A note on the 1,2-conjecture, in Ph.D. Thesis, 2009.
\bibitem{KKP} M. Kalkowski, M. Karo\'nski, F. Pfender, Vertex-coloring edge-weightings: Towards the 1-2-3 conjecture, {  Journal of Combinatorial Theory (B)} { 100} (2010), 347--349. 
 \bibitem{KLT}  M. Karo\'nski, T. \L uczak and A. Thomason,  Edge weights and vertex colours,  { J. Combinatorial Theory (B)}  { 91} (2004), 151--157.
 %\bibitem{P} F. Pfender, Total weight choosability in hypergraphs,  	arXiv:1312.6329 [math.CO].
 \bibitem{PW} J. Przyby{\l}o, M. Wo\'{z}niak, On a 1-2 conjecture, { Discrete Mathematics and Theoretical Computer Science} { 12} (2010), 101--108.
 \bibitem{PW2} J. Przyby{\l}o, M. Wo\'{z}niak, Total weight choosability of graphs, { Electronic J. Combinatorics} { 18} (2011), {\#}P112.
 \bibitem{T} C. Thomassen, The even cycle problem for directed graphs, { Journal of the AMS} { 5} (1992), 218--229.
%\bibitem{WY} T. Wang and Q. Yu, On vertex-coloring $13$-edge-weighting, {\em Front. Math. China} {\bf 3} (2008), 1-7.
\bibitem{WZ} T. Wong, X. Zhu, Total weight choosability of graphs, { J. Graph Theory} { 66} (2011), 198--212.
\bibitem{ZW} X. Zhu, T. Wong, Every graph is $(2,3)$-choosable, { Combinatorica} (2014), 1--7.
 \end{thebibliography}







\end{document}

